\documentclass[10pt,a4paper]{article}

% Web version of the CV: private contact details removed; DOIs added.
% Build: pdflatex cv-web.tex (twice), then copy cv-web.pdf to public/documents/cv-alexandros-kaminas.pdf
% ---------- Packages ----------
\usepackage[utf8]{inputenc}
\usepackage[T1]{fontenc}
\usepackage[default]{lato}
\usepackage{tgpagella}
\usepackage[a4paper,top=1.6cm,bottom=1.7cm,left=1.8cm,right=1.8cm]{geometry}
\usepackage[dvipsnames]{xcolor}
\usepackage[hidelinks]{hyperref}
\usepackage{titlesec}
\usepackage{enumitem}
\usepackage{fontawesome5}
\usepackage{fancyhdr}
\usepackage{lastpage}

% ---------- Colours ----------
\definecolor{ink}{HTML}{1B2430}
\definecolor{accent}{HTML}{0B5F80}
\definecolor{muted}{HTML}{5E6873}
\definecolor{rule}{HTML}{C9D3DB}

% ---------- Global formatting ----------
\setlength{\parindent}{0pt}
\setlength{\parskip}{0pt}
\setlength{\emergencystretch}{2em}
\raggedbottom
\urlstyle{same}
\color{ink}

\newcommand{\serif}{\fontfamily{qpl}\selectfont}

\pagestyle{fancy}
\fancyhf{}
\renewcommand{\headrulewidth}{0pt}
\fancyfoot[L]{\footnotesize\color{muted}Alexandros Kaminas \textperiodcentered{} Curriculum vitae}
\fancyfoot[R]{\footnotesize\color{muted}\thepage\,/\,\pageref{LastPage}}

\hypersetup{
  pdftitle={Curriculum Vitae -- Alexandros Kaminas},
  pdfauthor={Alexandros Kaminas},
  pdfsubject={Curriculum vitae (web version)},
  pdfkeywords={mathematical modelling, complex systems, eco-evolutionary dynamics, community ecology, sustainability}
}

% ---------- Section formatting ----------
\titleformat{\section}
  {\serif\Large\color{accent}}
  {}
  {0pt}
  {}
  [{\color{rule}\vspace{-3pt}\rule{\textwidth}{0.6pt}}]
\titlespacing*{\section}{0pt}{10pt}{5pt}

% ---------- Entry commands ----------
\newlength{\datew}
\setlength{\datew}{3.05cm}
\newlength{\bodyw}
\setlength{\bodyw}{\dimexpr\textwidth-\datew\relax}

% #1 date  #2 role (bold)  #3 organisation  #4 description (may be empty)
\newcommand{\entry}[4]{%
  \par\noindent
  \begin{minipage}[t]{\datew}\raggedright\footnotesize\leavevmode\color{muted}#1\end{minipage}%
  \begin{minipage}[t]{\bodyw}%
    \textbf{#2}\if\relax\detokenize{#3}\relax\else{\color{muted}\,\textperiodcentered\,}#3\fi
    \if\relax\detokenize{#4}\relax\else\par\vspace{1pt}{\small #4}\fi
  \end{minipage}\par
  \vspace{4pt}%
}

% Skills row. #1 label  #2 text
\newcommand{\skill}[2]{%
  \par\noindent
  \begin{minipage}[t]{\datew}\raggedright\small\leavevmode\color{muted}#1\end{minipage}%
  \begin{minipage}[t]{\bodyw}\small #2\end{minipage}\par
  \vspace{3pt}%
}

% Numbered publication with a date margin. #1 year  #2 text
\newcounter{pubno}
\newcommand{\pub}[2]{%
  \stepcounter{pubno}%
  \par\noindent
  \begin{minipage}[t]{\datew}\small\leavevmode\color{muted}[\arabic{pubno}]\hspace{0.6em}#1\end{minipage}%
  \begin{minipage}[t]{\bodyw}\small #2\end{minipage}\par
  \vspace{4pt}%
}

\newcommand{\me}{\textbf{Kaminas,~A.}}
\newcommand{\doi}[1]{\href{https://doi.org/#1}{\color{accent}doi:#1}}

% Key-figure cell. #1 figure  #2 label
\newcommand{\keyfig}[2]{%
  \begin{minipage}[t]{0.19\textwidth}\centering
    {\serif\huge\color{accent}#1}\\[1pt]
    {\footnotesize\color{muted}#2}
  \end{minipage}%
}

\begin{document}

% ================================================================
% Header
% ================================================================
{\serif\fontsize{30}{32}\selectfont Alexandros Kaminas}\par
\vspace{5pt}
{\large\color{accent}PhD researcher in theoretical and evolutionary ecology}\par
\vspace{2pt}
{\color{muted}University of Groningen, TR\^ES \textperiodcentered{} EVOLVE fellow (EU MSCA COFUND)}\par
\vspace{7pt}
{\small
\faEnvelope[regular]~\href{mailto:kaminas.alex@gmail.com}{kaminas.alex@gmail.com}\hspace{1.1em}
\faGlobe~\href{https://akaminas.github.io}{akaminas.github.io}\hspace{1.1em}
\faOrcid~\href{https://orcid.org/0000-0002-7928-651X}{0000-0002-7928-651X}\hspace{1.1em}
\faGithub~\href{https://github.com/akaminas}{akaminas}\hspace{1.1em}
\faGraduationCap~\href{https://scholar.google.com/citations?user=jyVZ4m4AAAAJ}{Google Scholar}
}\par
\vspace{6pt}
{\color{accent}\rule{\textwidth}{1.2pt}}\par
\vspace{8pt}

% ================================================================
% Profile
% ================================================================
I build mathematical and computational models of complex systems. My doctoral work at TR\^ES, under Rampal Etienne and Sander van Doorn, asks how isolation changes which ecological communities form and persist, and how they evolve. Before that I modelled regime shifts and functional diversity in Mediterranean fish communities at UMR MARBEC, spent five years monitoring species invasions through citizen science, and served as a meteorology and oceanography officer in the Hellenic Navy. Outside ecology I apply the same systems thinking to sustainability organisations and policy.

\vspace{9pt}
\noindent
\keyfig{10+}{peer-reviewed articles}\hfill
\keyfig{3}{journals reviewed for}\hfill
\keyfig{4}{fellowships and awards}\hfill
\keyfig{3}{countries of research}\hfill
\keyfig{3}{working languages}\par
\vspace{8pt}

{\small\textbf{Research interests}\quad
Eco-evolutionary dynamics \textperiodcentered{}
community assembly \textperiodcentered{}
isolation and immigration \textperiodcentered{}
diversification and extinction \textperiodcentered{}
regime shifts \textperiodcentered{}
functional diversity \textperiodcentered{}
biological invasions \textperiodcentered{}
sustainability systems}\par

% ================================================================
% Education
% ================================================================
\section*{Education}

\entry{Oct 2025 -- present}
{PhD, The Evolutionary Consequences of Isolation}
{University of Groningen, Netherlands}
{Theoretical Research in Evolutionary Life Sciences (TR\^ES). Working under Rampal S.~Etienne (Theoretical and Evolutionary Community Ecology) and G.~Sander van Doorn (Evolutionary Systems Biology). An eco-evolutionary modelling framework linking isolation, immigration, ecological interactions, diversification and extinction; stochastic simulation in C++.}

\entry{Sep 2022 -- Jun 2024}
{MSc, Applied Ecological and Evolutionary Sciences}
{University of Montpellier, France}
{Graduated 2nd in class. IDIL international scholarship. Thesis research at UMR MARBEC on regime shifts in Mediterranean communities.}

\entry{Sep 2017 -- Dec 2021}
{BSc, Marine Science}
{University of the Aegean, Greece}
{Graduated 1st in class (top 1\%) and delivered the valedictory address.}

% ================================================================
% Fellowships & awards
% ================================================================
\section*{Fellowships and Awards}

\entry{Oct 2025}{EVOLVE Fellowship}{EU EVOLVE MSCA COFUND}{Funds the doctoral project at the University of Groningen.}
\entry{Aug 2022}{IDIL International Student Scholarship}{University of Montpellier}{}
\entry{Jan 2022}{Academic Excellence Award}{Greek State Scholarships Foundation (IKY)}{}
\entry{Dec 2021}{Valedictorian}{University of the Aegean}{}

% ================================================================
% Research experience
% ================================================================
\section*{Research Experience}

\entry{Jan -- Jun 2024}
{Research Intern}
{UMR MARBEC, Montpellier, France}
{Regime shifts and community reorganisation of Mediterranean fish and invertebrates under climate change and fishing, using catastrophe theory, stochastic modelling and multi-effects models.}

\entry{Jan -- Jun 2023}
{Research Intern}
{UMR MARBEC, S\`ete, France}
{Functional diversity of fish communities with $n$-dimensional hypervolumes; phylogenetic inference and trait imputation.}

\entry{Sep -- Dec 2021}
{Erasmus Research Intern}
{UMR MARBEC, S\`ete, France}
{Statistical analysis of functional diversity under global change.}

\entry{Jun -- Sep 2020}
{Research Intern}
{Hellenic Centre for Marine Research (HCMR), Athens}
{Taxonomy, stomach-content analysis and diet indices of commercial fish species; led to a co-authored article in \emph{Animals} [3].}

\entry{2018 -- 2023}
{Volunteer Research Fellow}
{iSea Environmental Organisation, Greece}
{Monitoring of invasive species and elasmobranchs; co-authored peer-reviewed articles on species records and range expansions.}

\entry{2021 -- present}
{Working Group Member}
{WGCOMEDA}
{Comparative analysis of Mediterranean and Atlantic ecosystems, biodiversity and resilience.}

% ================================================================
% Publications
% ================================================================
\section*{Selected Publications}

\pub{2026}{Alevizos, A., \me, Vrouvaki, C., \& Papandreou, A. Assessing organizational commitment to sustainable development and circular economy principles among Greek NGOs: an analysis using organizational culture, institutional isomorphism, and systems frameworks. \emph{Circular Economy and Sustainability}, 6(3), 204. \doi{10.1007/s43615-026-00768-8}}

\pub{2023}{Ragkousis, M., Zenetos, A., Souissi, J.\,B., Hoffman, R., Ghanem, R., Ta\c{s}k\i n, E., \ldots{} \me{} Unpublished Mediterranean and Black Sea records of marine alien, cryptogenic, and neonative species. \emph{BioInvasions Records}, 12(2), 339--369. \doi{10.3391/bir.2023.12.2.01}}

\pub{2023}{Mina, A., Mytilineou, C., \me, Rekleiti, A., Siapatis, A., \& Anastasopoulou, A. Feeding habits of \emph{Dentex maroccanus} and the effect of body size. \emph{Animals}, 13(5), 939. \doi{10.3390/ani13050939}}

\pub{2023}{Minasidis, V., Doumpas, N., Giovos, I., Kleitou, P., \me, \& Moutopoulos, D.\,K. Assessing consumer attitude towards marine non-indigenous fish species: a case study from Greece. \emph{Thalassas}, 39(1), 35--53. \doi{10.1007/s41208-022-00486-6}}

\pub{2022}{\me, Shokouros-Oskarsson, M., Minasidis, V., Langeneck, J., Kleitou, P., Tiralongo, F., \& Crocetta, F. Filling gaps via citizen science: \emph{Phyllorhiza punctata} von Lendenfeld, 1884 (Cnidaria: Scyphozoa: Mastigiidae) in Cyprus. \emph{BioInvasions Records}, 11(3), 667--675. \doi{10.3391/bir.2022.11.3.09}}

\pub{2022}{Langeneck, J., Minasidis, V., Doumpas, N., Giovos, I., \me, Kleitou, P., Tiralongo, F., \& Crocetta, F. Citizen science helps in tracking the range expansions of non-indigenous and neo-native species in Greece and Cyprus. \emph{Journal of Marine Science and Engineering}, 10(2), 256. \doi{10.3390/jmse10020256}}

\pub{2021}{Giovos, I., Serena, F., Katsada, D., \ldots{} \me, \ldots{} \& Kleitou, P. Integrating literature, biodiversity databases, and citizen-science to reconstruct the checklist of chondrichthyans in Cyprus. \emph{Fishes}, 6(3), 24. \doi{10.3390/fishes6030024}}

\pub{2020}{Giovos, I., Tiralongo, F., Langeneck, J., \me, Kleitou, P., Crocetta, F., \& Doumpas, N. First record of the Atlantic spadefish \emph{Chaetodipterus faber} (Broussonet, 1782) in the Mediterranean Sea: is it a new aquarium release? \emph{BioInvasions Records}, 9(1), 89--95. \doi{10.3391/bir.2020.9.1.12}}

% ================================================================
% Presentations
% ================================================================
\section*{Conference Presentations}

\entry{Jun 2026}
{Oral \textperiodcentered{} NVTB Annual Schoorl Meeting}
{Schoorl, Netherlands}
{\emph{Isolation as a dynamic driver of eco-evolutionary change}. With G.\,S.~van Doorn and R.\,S.~Etienne.}

\entry{Jun 2024}
{Oral \textperiodcentered{} 16th AFH Colloquium}
{S\`ete, France}
{Effects of climate and fishing on Mediterranean communities, analysed with catastrophe theory. The only MSc project selected for an oral presentation.}

\entry{2021}
{Two posters \textperiodcentered{} Hydromedit}
{Mytilene, Greece}
{Rare-species monitoring through citizen science; diet composition of \emph{Dentex maroccanus}.}

% ================================================================
% Supervision, review & service
% ================================================================
\section*{Supervision, Peer Review and Service}

\entry{2026}
{MSc Research Supervisor}
{University of Groningen, TR\^ES}
{Supervised the MSc project \emph{Modelling Trait Evolution in Mutualistically Interacting Networks}.}

\entry{Nov 2026 -- present}{Peer Reviewer}{\emph{Ecological Informatics}}{Manuscripts on complex systems models.}
\entry{Sep 2026 -- present}{Peer Reviewer}{\emph{Ecological Modelling}}{Manuscripts on ecological and eco-evolutionary models.}
\entry{Jan 2025 -- present}{Peer Reviewer}{\emph{Heliyon}}{Manuscripts on environmental modelling and mathematics.}

\entry{Jun 2026}
{Financial Auditor}
{Netherlands Society for Theoretical Biology (NVTB)}
{Audited the society's finances as part of its financial oversight.}

\entry{Jan 2024 -- present}
{International Student Ambassador}
{Universit\'e de Montpellier}
{International student outreach and public speaking.}

% ================================================================
% Software & teaching
% ================================================================
\section*{Scientific Software and Teaching}

\entry{2026}
{Vellum}
{\href{https://akaminas.github.io/vellum/}{\color{accent}akaminas.github.io/vellum}}
{Author. A visual style for scientific figures: a computed colour palette, a fixed figure anatomy and style files for Python, R, \LaTeX{} and the web.}

\entry{2026}
{C++ Atlas}
{\href{https://akaminas.github.io/cppverse/}{\color{accent}akaminas.github.io/cppverse}}
{Author. An independent course that teaches C++ to biologists as a route to scientific modelling.}

\entry{2026}
{Tessera}
{\href{https://akaminas.github.io/Tessera/}{\color{accent}akaminas.github.io/Tessera}}
{Author. A project-management web app with work packages, budgets and reporting for funded multi-partner projects (React, Supabase).}

\entry{Jan 2018 -- Jun 2021}
{Academic Tutor}
{Mytilene, Greece}
{Biology, geology, physics and chemistry; coursework, thesis preparation and study skills.}

% ================================================================
% Sustainability & science policy
% ================================================================
\section*{Sustainability and Science Policy}

\entry{May 2025 -- present}
{Network Coordinator}
{SDSN Youth Greece}
{Coordinate collaborations and projects on the Sustainable Development Goals. Previously Sustainability Research Fellow (Nov 2023 -- May 2025), managing projects and co-authoring SDG publications.}

\entry{Jan 2025 -- present}
{Research Associate}
{Political Economy of Sustainable Development Lab, University of Athens}
{Climate-change reporting, data visualisation and science communication; systems analysis of NGOs [1].}

\entry{2024 -- 2025}
{Project Contributor}
{SDG4U, EU-funded (Erasmus+) higher-education project}
{Research, writing and dissemination. \href{https://www.sdg4uproject.eu/team/}{\color{accent}Project team page}.}

% ================================================================
% Citizen science & public engagement
% ================================================================
\section*{Citizen Science and Public Engagement}

\entry{2019 -- 2023}
{Marine Alien Species Population Density in Greece and Cyprus}
{}
{Citizen-science monitoring of alien-species population density.}

\entry{2018 -- 2023}
{Is it Alien to you? Share it!!!}
{}
{Public engagement and alien-species monitoring under a research permit from the Greek Ministry of Environment.}

\entry{2018 -- 2023}
{MECO Project}
{}
{Collaborative network recording elasmobranch sightings across the Mediterranean.}

\entry{2018 -- 2019}
{Plastic Litter Project}
{Marine Remote Sensing Club, Lesvos}
{Detection of marine plastic litter from satellite imagery.}

% ================================================================
% Professional experience
% ================================================================
\section*{Professional Experience}

\entry{Sep 2024 -- Sep 2025}
{Officer, Hellenic Navy}
{Athens and Poros, Greece}
{\emph{Meteorology and Oceanography Officer (METOC)}, Hydrographic Service (Oct 2024 -- Sep 2025): worked with oceanographic models, issued forecasts and processed acoustic data for submarine surveys.\\
\emph{Administration Officer}: managed databases covering more than 600 trainees, modernised administrative systems with Python, and handled sensitive headquarters data.}

% ================================================================
% Skills
% ================================================================
\section*{Skills and Languages}

\skill{Programming}{C++ (advanced), Python (advanced), R (advanced), MATLAB (basic)}
\skill{Methods}{Stochastic simulation, eco-evolutionary models, catastrophe theory, trait-space (hypervolume) analysis, phylogenetic inference and trait imputation}
\skill{Spatial \& imaging}{QGIS, ArcGIS, Image-Pro V9}
\skill{Laboratory}{Sample treatment for pollution analysis; image-based identification; stomach-content analysis}
\skill{Languages}{Greek (native), English (C2), French (C1)}

\end{document}
